\section{Application patterns and repairs of packet difficulties}

The framework is intended as a shared substrate.  This section states what it
does and does not resolve in the principal domains represented in the
repository.

\subsection{Algebraic equations and Puiseux expansions}

For a polynomial with exact rational or exact algebraic coefficient codes,
finite algebraic operations, resultants, discriminants, square-free
factorization certificates, and coefficient identities are decidable finite
data.  A root theorem should first output a finite list of rational disks with:
coverage of all roots, multiplicity counts, and a refinement relation between
precisions.  A separation certificate collapses an isolated disk to one root
name.

For arbitrary Cauchy-presented coefficients, exact Puiseux coefficient
selection is too strong.  The repository gives an explicit fixed-family
halting reduction using
\[
       P_a(X,Y)=Y^2-X^2-a.
\]
The repair is one of:
\begin{enumerate}[label=(\alph*)]
\item exact algebraic coefficient presentations with decidable zero and
      certified splitting;
\item a finite family of truncated approximate branches with residual and
      coverage bounds;
\item a theorem conditional on a simple-root or discriminant gap.
\end{enumerate}
The contract-DAG compiler handles the coefficient recursion once the branch
guard is legitimate; it does not create the missing exact equality oracle.

\subsection{Euler--Rodrigues, eigenvectors, and projective geometry}

Rodrigues' formula for a supplied axis and angle is a polynomial/trigonometric
error-contract problem.  Exact Euler-axis selection for an arbitrary
represented rotation is discontinuous at the identity and can encode LLPO.
The stable outputs are:
\begin{itemize}[leftmargin=1.6em]
\item a unit quaternion modulo \(q\sim-q\);
\item a projective axis or finite approximate axis set;
\item an exact axis under a supplied lower bound separating the rotation from
      the identity;
\item rational exact axes for rational proper orthogonal matrices.
\end{itemize}
The same projective-first discipline applies to eigenvectors.  Use a
projector, kernel certificate, or finite set of charts; normalize only under a
nonzero coordinate or singular-value gap.

Many Euclidean and projective geometry theorems become polynomial identities
after homogeneous coordinates are introduced.  Denominators should not be
cancelled silently.  Either clear them and retain nondegeneracy guards, or
work entirely with cross-products and determinants.  A finite incidence
certificate then fits directly into the contract calculus.

\subsection{Analysis and special functions}

A valid analysis node has:
\[
   \text{approximant},\quad
   \text{cutoff},\quad
   \text{tail/error verifier},\quad
   \text{precision demand}.
\]
For a series, the producer supplies \(N(k,d)\) and a rational proof that the
tail after \(N\) is at most \(2^{-k}\).  For an integral, it supplies a mesh
and local oscillation modulus.  For analytic continuation, it supplies a
finite chain of overlapping disks and derivative bounds.  The compiler adds
the budgets and verifies the finite chain.

This resolves a common formal difficulty in trigonometric, Gamma, zeta, and
residue packets: one never inducts on a completed real identity.  One proves
a rational inequality at arbitrary requested precision.  It does not, by
itself, prove a missing global zero-free region or manufacture an effective
constant from an ineffective theorem.

\subsection{Probability}

Finite probability laws are rational mass arrays with a Boolean check that
masses are nonnegative and sum to one.  Convolution, moments, and
characteristic-function samples are exact finite sums.  A constructive CLT
packet should separate:
\begin{enumerate}[label=(\roman*)]
\item the exact finite convolution engine;
\item moment and variance guards;
\item a finite-grid comparison or characteristic-function approximation;
\item the analytic smoothing and tail estimates with explicit constants.
\end{enumerate}
The first three are generic contract nodes.  The fourth is the substantive
analytic producer.

\subsection{Topology}

For Jordan-curve and invariance-of-dimension arguments, use finite grids,
rational boxes, adjacency graphs, parity or degree certificates, and explicit
separation margins.  A finite grid calculation at a few scales is evidence,
not a proof for every scale.  The reusable theorem must show that the
quantifier-free grid invariant is preserved under refinement.

Degenerate boundary cases should be handled by overlap: several adjacent
cells may be valid.  A positive distance from the boundary collapses the
finite relation to one cell.  This is the topological instance of the same
branching theorem used for roots and axes.

\subsection{Diophantine geometry}

Runge-type arguments often do produce explicit polynomial identities and
height inequalities.  These are ideal evidence-indexed nodes: verify the
identity, derive the bound, then enumerate.  The general Siegel and Faltings
packets correctly identify the missing general height producer.  The
framework contributes:
\begin{itemize}[leftmargin=1.6em]
\item exact encodings of rational and integral points;
\item complete bounded enumeration;
\item conversion between a verified bound and a complete list;
\item finite-cover and finite-algebra verifiers;
\item a precise interface for any future height or covering theorem.
\end{itemize}
A cardinality theorem, even a strong one, cannot be plugged into the height
input slot.

\subsection{Analytic number theory}

Prime-distribution statements must expose all dependencies:
modulus range, exceptional characters, zero-free constants, conductor cutoff,
density estimate, and finite small cases.  A certificate verifier can check
the finite exceptional list and derive a claimed numerical range from
explicit inequalities.  The quantifier order is crucial:
\[
  \exists C\;\forall x\ge X(C)\;\Phi(C,x)
\]
is useful only when \(C\) is produced or supplied.  A proof of mere existence
in a stronger classical theory does not automatically yield a PR constant in
iPRA.  The framework therefore records ``analytic constant producer'' as a
first-class dependency rather than burying it in a phrase such as ``take
\(x\) sufficiently large.''

\subsection{Summary matrix}

\begin{longtable}{p{.22\linewidth}p{.34\linewidth}p{.34\linewidth}}
\toprule
\textbf{Domain} & \textbf{Uniform output supported} &
\textbf{Additional evidence for exact choice}\\
\midrule
\endhead
Polynomial roots & finite disk cover, residuals, multiplicities &
root separation or exact algebraic splitting\\
Puiseux & finite approximate truncations &
exact coefficient field and branch isolation\\
Rotations & quaternion modulo sign, approximate axis set &
distance from identity / nonzero vector part\\
Eigenvectors & projective line or projector &
spectral gap and normalization chart\\
Limits and integrals & rational approximants with moduli &
none beyond the quantitative hypotheses\\
Topology & finite grid/parity relation &
boundary separation for unique cell/branch\\
Diophantine points & complete list below supplied height &
height or covering certificate producer\\
Prime distribution & verifier for explicit constants/ranges &
effective constant and exceptional-set producer\\
\bottomrule
\end{longtable}
